% concatenated from 2parametersversion.tex
\documentclass[12pt]{article}


% Packages used
\usepackage{amsfonts,amsmath,amsthm,bm,bbm, thmtools, thm-restate}
\usepackage{latexsym,color,epsfig,mathrsfs,enumerate}
\usepackage[toc,page]{appendix}
\usepackage{fancyhdr}
\usepackage[colorlinks=true, allcolors=red]{hyperref}
\usepackage{indentfirst}
\usepackage[maxbibnames=9]{biblatex}
\usepackage{comment}
\usepackage{xcolor}
\usepackage[shortlabels]{enumitem}

\usepackage{graphicx}
\usepackage{pgfplots}
\pgfplotsset{compat=1.15}
\usetikzlibrary{arrows}

\addbibresource{bibliography.bib}
\nocite{*}



% Dimensions of page
\setlength{\textheight}{22.5cm} 
\setlength{\textwidth}{6.7in}
\setlength{\topmargin}{0pt} 
\setlength{\evensidemargin}{1pt}
\setlength{\oddsidemargin}{1pt} 
\setlength{\headsep}{10pt}
\setlength{\parskip}{2mm} 
\setlength{\parindent}{5mm}
\renewcommand{\baselinestretch}{1.1}


% Environments
\newtheorem{question}{Question}
\newtheorem{thm}{Theorem}[section]
\newtheorem{dfn}[thm]{Definition}
\newtheorem{lem}[thm]{Lemma}
\newtheorem{cor}[thm]{Corollary}
\newtheorem{prop}[thm]{Proposition}
\newtheorem{rmk}[thm]{Remark}
\newtheorem{obs}[thm]{Observation}
\newtheorem{conjecture}[thm]{Conjecture}
\newtheorem{exam}{Example}


% Commands with argument
\newcommand{\mc}[1]{\mathcal{#1}}
\newcommand{\mb}[1]{\mathbb{#1}}
\newcommand{\nib}[1]{\noindent {\bf #1}}
\newcommand{\nim}[1]{\noindent {\em #1}}
\newcommand{\brac}[1]{\left( #1 \right)}
\newcommand{\bracc}[1]{\left\{ #1 \right\}}
\newcommand{\bsize}[1]{\left| #1 \right|}
\newcommand{\brak}[1]{\left[ #1 \right]}
\newcommand{\bangle}[1]{\left\langle #1 \right\rangle}
\newcommand{\bfl}[1]{\left\lfloor #1 \right\rfloor}
\newcommand{\bcl}[1]{\left\lceil #1 \right\rceil}
\newcommand{\bin}[2]{\binom {#1}{#2}}
\newcommand{\sbin}[2]{\binom {[#1]}{#2}}
\newcommand{\bing}[2]{\binom {#1}{\geq #2}}
\newcommand{\sbing}[2]{\binom {[#1]}{\geq #2}}
\newcommand{\binl}[2]{\binom {#1}{\leq #2}}
\newcommand{\sbinl}[2]{\binom {[#1]}{\leq #2}}
\newcommand{\laur}[1]{\textbf{\underline{#1}}}
\newcommand{\bad}[1]{\operatorname{bad}_{\mathcal{D}}({#1})}
\newcommand{\baddd}[4]{\operatorname{bad}_{\mathcal{#1}_{#2}}^{#3}({#4})}
\newcommand{\proj}[3]{\operatorname{proj}_{#1}({\bf {#2}}_{#3})}
\newcommand{\diver}[1]{\operatorname{div}(#1)}
\newcommand{\dd}[2]{\operatorname{deg-diff}^{#2}({#1})}
\newcommand{\homm}[1]{\operatorname{hom}(#1)}
%\newcommand{\laurp}[1]{\textbf{\underline{#1^\prime}}}
% General structure 
% \newcommand{name}[num]{definition}
% \newcommand{\FunctionName}[2]{funct(#1 , #2)}

\usepackage{combelow}% provides \cb to place comma below character

\usepackage{newunicodechar}
\newunicodechar{ș}{\cb{s}}
\newunicodechar{ț}{\cb{t}}

% Commands without argument
\newcommand{\sub}{\subset}
\newcommand{\subn}{\subsetneq}
\newcommand{\sups}{\supseteq}
\newcommand{\ex}{\mbox{ex}}
\newcommand{\ra}{\rightarrow}
\newcommand{\Ra}{\Rightarrow}
\newcommand{\sm}{\setminus}
\newcommand{\es}{\emptyset}
\newcommand{\pl}{\partial}
\newcommand{\ov}{\overline}
\newcommand{\ul}{\underline}
\newcommand{\wt}{\widetilde}
\newcommand{\eps}{\varepsilon}


% Character Abbreviations
\newcommand{\aA}{\alpha}
\newcommand{\bB}{\beta}
\newcommand{\gG}{\gamma}
\newcommand{\dD}{\delta}
\newcommand{\iI}{\iota}
\newcommand{\kK}{\kappa}
\newcommand{\zZ}{\zeta}
\newcommand{\lL}{\lambda}
\newcommand{\tT}{\theta}
\newcommand{\sS}{\sigma}
\newcommand{\oO}{\omega}
% Larger Character Abbreviations
\newcommand{\GG}{\Gamma}
\newcommand{\DD}{\Delta}
\newcommand{\OO}{\Omega}
\newcommand{\TT}{\Theta}
\newcommand{\Ss}{\Sigma}
\newcommand{\LL}{\Lambda}





\newcommand{\A}{\mathcal{A}}
\newcommand{\B}{\mathcal{B}}
\newcommand{\C}{\mathcal{C}}
\newcommand{\Pn}{{\mathcal P}[n]}
\newcommand{\R}{\mathbb R}
\newcommand{\N}{\mathbb N}
\newcommand{\bP}{\mathbb P}
\newcommand{\bE}{\mathbb E}

\title{Many Antipodal Pairs Force Many Neighboring Pairs}
\date{}


\author{G\'abor Dam\'asdi\thanks{HUN-REN Alfr\'ed R\'enyi Institute of Mathematics, ELTE Eötvös Loránd University, Budapest. Email: \texttt{damasdigabor@caesar.elte.hu}} \and Lauren\cb{t}iu Ploscaru\thanks{HUN-REN Alfr\'ed R\'enyi Institute of Mathematics, Budapest. Email: \texttt{laurentiuploscaru@renyi.hu}.
}}



\begin{document}

\maketitle

\begin{abstract}
	  Let $X=\{x_1,\dots,x_n\}\subset \R^2$ be a finite set of points of diameter at most $1$. It is natural to expect that if many pairs $(x_i,x_j)$ lie at distance close to $1$ from each other, then some clustering phenomenon must occur, implying that a significant number of these pairs are also very close to each other. %For $0<\eps<1$, we call a pair $(x_i,x_j)$ $\varepsilon$-\emph{antipodal} if $\|x_i-x_j\|\ge 1-\eps$ and $\varepsilon$-\emph{neighboring} if $\|x_i-x_j\|\le \eps$.
      We prove that there exists a universal constant $c>0$ such that for all $0<\eps<1$, whenever $n$ is large enough, we have: 
\[
\big|\{(i,j):\|x_i-x_j\|\le \eps\}\big|
\ge c\cdot \eps^{1/2}\cdot
\big|\{(i,j):\|x_i-x_j\|\ge 1-\eps\}\big|.
\]

This confirms a recent conjecture of Steinerberger, who asked whether the $\eps^{1/2}$ ratio is the best possible. We also study a two-parameter version of Steinerberger's question by considering the number of pairs at distance at most $\eps_1$ and at distance at least $1-\eps_2$. We show that in this case the optimal ratio is $\eps_1^2\cdot\eps_2^{-3/2}$. The proof proceeds by introducing an auxiliary graph associated with the set $X$ and reducing the problem to bounding the largest eigenvalue of its adjacency matrix. Our main result is the outcome of human--AI interactions using ChatGPT 5.4. 
      
\end{abstract}


\section{Introduction}

 Consider a finite set of points  $X=\{x_1,\dots,x_n\}\subset \R^2$ of diameter at most 1. What can we say about the number of pairs of points that are within distance $\varepsilon$ from each other? The area of the convex hull is bounded, hence we always have  $\Omega(n^2\varepsilon^{2})$ such pairs, and this cannot be improved in general. In this paper we investigate how this behavior changes if we assume that there are many pairs of points in $X$ at distance at least $1-\varepsilon$. 

 For $0<\eps<1$, we call a pair $(x_i,x_j)$ $\varepsilon$-\emph{antipodal} if $\|x_i-x_j\|\ge 1-\eps$ and $\varepsilon$-\emph{neighboring} if $\|x_i-x_j\|\le \eps$.  Let $\mathcal{A}_\eps(X)$ and $\mathcal{N}_\eps(X)$ denote the set of antipodal pairs and the set of neighboring pairs, respectively. More precisely, we define:
\begin{align*}
     \mathcal{A}_\eps(X):=&\{\{x,y\}\subset X: \|x-y\|\geq 1-\eps\}, \\
     \vspace{4pt}\mathcal{N}_\eps(X):=&\{\{x,y\}\subset X: \|x-y\|\leq \eps\}. 
\end{align*}   
  It is natural to expect that if we have many antipodal pairs, then some of the points are forced to be close to the boundary and the number of neighboring pairs increases. Furthermore, as we increase the number of antipodal pairs, we can also expect some clustering of the points. Hence, we are interested in how small the ratio $R_\eps(X):=|\mathcal{N}_\eps(X)|/|\mathcal{A}_\eps(X)|$  can be. This line of research was recently initiated by Steinerberger \cite{steinerberger2025antipodesimpliesneighbors}. He showed the following: %theorem. 

\begin{thm}[Steinerberger \cite{steinerberger2025antipodesimpliesneighbors}]\label{thm:steinerberger}
There are constants $c_1,c_2>0$ such that for all $0<\varepsilon <1$ and any finite $X\subset \mathbb{R}^2$ of diameter at most $1$ with $|X|\ge c_1\varepsilon^{-2}$, we have $R_\eps(X)\ge c_2\sqrt[4]{\eps^3/\log(\eps^{-1})}$, i.e.
    \[
\big|\{(i,j):\|x_i-x_j\|\le \eps\}\big|
\ge c_2\cdot \frac{\eps^{3/4}}{(\log\eps^{-1})^{1/4}}\cdot
\big|\{(i,j):\|x_i-x_j\|\ge 1-\eps\}\big|.
\]

\end{thm}

Steinerberger provided a number of examples where the ratio is of order $\eps^{1/2}$, one of which is a regular polygon on sufficiently many vertices, and asked whether this threshold is optimal. Very recently, Korsky \cite{korsky2026optimalspectralboundsantipodal} showed that this is achievable up to a polylogarithmic factor. 

\begin{thm}[Korsky \cite{korsky2026optimalspectralboundsantipodal}]\label{thm:korsky}
There are universal constants $c_1,c_2>0$ such that for any $0<\varepsilon<1$  and  any finite $X\subset \mathbb{R}^2$ of diameter at most $1$ with $|X|\ge c_1\varepsilon^{-2}$, we have ${R_\eps(X)\ge c_2\sqrt{\eps/\log(\eps^{-1})}}$, i.e.
    \[
\big|\{(i,j):\|x_i-x_j\|\le \eps\}\big|
\ge c_2\cdot \frac{\eps^{1/2}}{(\log\eps^{-1})^{1/2}}\cdot
\big|\{(i,j):\|x_i-x_j\|\ge 1-\eps\}\big|.
\]

\end{thm}

We note that in both papers the authors stated the result without assuming any lower bound on the number of points in $X$. Technically, their statements are incorrect. In fact, the condition $|X|=\Omega(\varepsilon^{-2})$ is necessary since one can pack $\Omega(\varepsilon^{-2})$ axis-parallel squares of side-length $1.01\cdot\eps$ into the disk of diameter 1 and take $X$ to be the set of their centers. In this case the ratio $R_\eps(X)$ is 0 since $|\mathcal{N}_\eps(X)|=0$, and no $c_2>0$ works. We have  presented Theorem \ref{thm:steinerberger} and Theorem~\ref{thm:korsky} in their corrected form. 

In this paper we confirm Steinerberger's conjecture. 

\begin{thm}\label{thm:main_single}
There are universal constants $c_1,c_2>0$ such that for any $0<\varepsilon<1$  and  any finite $X\subset \mathbb{R}^2$ of diameter at most $1$ with $|X|\ge c_1\varepsilon^{-2}$ we have $R_\eps(X)\ge c_2\eps^{1/2}$, i.e.
    \[
\big|\{(i,j):\|x_i-x_j\|\le \eps\}\big|
\ge c_2\cdot \eps^{1/2}\cdot
\big|\{(i,j):\|x_i-x_j\|\ge 1-\eps\}\big|.
\]

\end{thm}

Theorem \ref{thm:main_single} is a consequence of a more general result. We study the relationship between $|\mathcal{N}_{\eps_1}(X)|$ and $|\mathcal{A}_{\eps_2}(X)|$ for parameters $0<\eps_1,\eps_2<1$, focusing on the minimum possible value of their ratio. Our main result determines this quantity up to a constant factor.     

\begin{thm}\label{thm:main}
There are universal constants $c_1,c_2>0$ such that for any $0<\varepsilon_1,\varepsilon_2<1$ and any finite set $X\subset \mathbb{R}^2$ of diameter at most $1$ with $|X|\ge c_1\varepsilon_1^{-2}$ the following holds: 
\[
\dfrac{|\mathcal{N}_{\eps_1}(X)|}{|\mathcal{A}_{\eps_2}(X)|} \ge
\begin{cases}
c_2\cdot \eps_1^2\cdot \eps_2^{-3/2}, & \eps_1 \le \eps_2; \\
 c_2\cdot \eps_1^{1/2}, & \eps_1 > \eps_2.
\end{cases}
\]
\end{thm}

Theorem~\ref{thm:main}, together with the constructions in Section~\ref{sec:examples}, determines the optimal order of magnitude throughout the whole parameter range. In particular, it yields a tight resolution to Steinerberger's question, as it implies that $R_\eps(X)\geq c_2\eps^{1/2}$. Throughout the paper we focus on the two-parameter version. We will see examples showing that Theorem \ref{thm:main} is optimal, and we will discuss why a transition in behavior occurs at  $\eps_1=\eps_2$.

As far as the authors are aware, this is the first comprehensive result in this direction of study. There are a couple of results whose aim is to maximize the average distance between two points within a set of bounded diameter. For example, Witsenhausen \cite{Witsenhausen1974} studied the problem of maximizing $\sum_{i,j}\|x_i-x_j\|^2$ for a finite set $X:=\{x_1,x_2,\ldots,x_n\}\subset \mathbb{R}^2$ of diameter $1$, while Pillichshammer \cite{Pillichshammer2001, Pillichshammer2000O} analyzed $\sum_{i,j}\|x_i-x_j\|$. Theorem \ref{thm:main} relates to these results in the sense that having many pairs with maximum contribution (at least $1-\eps$) forces the existence of many pairs with a very small contribution (less than $\eps$). 

The proof of Theorem \ref{thm:main} was completed in close collaboration with the artificial intelligence ChatGPT 5.4.  
The spirit of our argument comes fairly close to the proofs of Theorem \ref{thm:steinerberger} and Theorem \ref{thm:korsky}. Following the ideas of Steinerberger, we begin by observing that only the points in $X$ that are within $\eps_2$ distance from the boundary of $\text{conv}(X)$ can contribute to $\eps_2$-antipodal pairs. Then we simplify the problem by covering the region near the boundary of $\text{conv}(X)$ with boxes of diameter less than $\eps_1$. When $\eps_1\leq \eps_2$, a natural choice for these boxes is given by squares of side length $\eps_1/2$. On the other hand when $\eps_1>\eps_2$, and especially when $\eps_2=o(\eps_1)$, these squares would cover a much larger area than needed, and a better choice is given by rectangles of dimension $\eps_1/2\times \eps_2/2$. 
This also explains the dichotomy in the conclusion of Theorem \ref{thm:main}. Next, we define a graph $G_X$, called the \emph{box graph}, on these boxes by connecting pairs that are at least $1-\eps_2$ far from each other. The problem reduces to bounding the largest eigenvalue of the adjacency matrix of the box graph. Using a key geometric lemma, we show that $G_X$ satisfies a natural sparsity condition. Roughly speaking, there is a universal constant $c>0$ such  that any vertex has at most $c|V(G_X)|\cdot (\eps_2/\eps_1)^{2} t^{-1}$ neighbors of degree at least $t$ for any integer $t>0$. This property allows us to derive the optimal bound on the eigenvalue of the box graph.  

After introducing the necessary metric and linear algebraic tools in Section \ref{sec:linalg}, we perform the reduction to the box graph in Section \ref{sec:reduction}. In Section \ref{sec:geom} we present Lemma \ref{quadrilateral}, which encompasses the geometry of the box graph and shows how it implies the sparsity condition (Lemma \ref{lem:tails}). This allows us to finish the proof of Theorem \ref{thm:main}. Next, in Section \ref{sec:examples} we discuss various extremal examples which demonstrate the sharpness of our bounds. Finally, in Section \ref{sec:conclude}, we discuss some open problems and provide additional insight into the selection of test vectors we have investigated for the Collatz-Wielandt bound before reaching the effective choice from Proposition \ref{prop:sqrt-control}.







\section{Linear algebra and metric prerequisites}\label{sec:linalg}

In this short section, we recall a couple of useful results from linear algebra. The terminology we use is quite standard and our presentation is largely self-contained; however, the reader can consult \cite{booklinalg} for a more detailed discussion.

%\begin{dfn}
 %   Given a $n$-vertex (simple) graph $G$ with vertices $v_1,v_2,\ldots, v_n$, its adjacency matrix is a $n\times n$ binary symmetric matrix $A:=(a_{ij})_{1\leq i,j,\leq n}$ whose $1$-entries are given by the edges of $G$, i.e. $a_{ij}=a_{ji}=1$ if and only if there is an edge in $G$ between $v_i$ and $v_j$. 
%\end{dfn}


The results presented below could be generalized to a broader context, but for our purposes, it is enough to restrict ourselves to real-valued symmetric matrices. What is essential in this setting is that the eigenvalues of such matrices are real-valued, and the largest, denoted by $\lambda_1(A)$ for a matrix $A$, satisfies $\lambda_1(A)=\underset{\mathbf{x}\in \mathbb{R}^n\setminus\{0\}}{\max}\dfrac{\langle \mathbf{x},A\mathbf{x} \rangle}{\langle \mathbf{x},\mathbf{x} \rangle}$. 

When the symmetric matrix $A$ also has only non-negative entries, by the Perron-Frobenius theorem there is an eigenvector $\mathbf{x}\in \mathbb{R}^n$ for $\lambda_1(A)$ whose entries are all non-negative. This observation is crucial to proving the following useful proposition. 

\begin{prop}[Collatz--Wielandt bound]\label{vector-test-criterion}
 Let $A$ be a real symmetric $n\times n$ non-zero matrix with non-negative entries. Suppose there is a vector $\mathbf{x}\in \mathbb{R}^n_{> 0}$ and $\lambda>0$ such that $(A\mathbf{x})_i\leq \lambda x_i$ for each coordinate $i\in [n]$. Then $$\lambda\geq \lambda_1(A)>0.$$
\end{prop}

\begin{proof}
  Recall that $\lambda_1(A)=\max_{||v||=1}\langle v,Av\rangle$, so clearly $\lambda_1(A)>0$. Moving forward, pick an eigenvector $v$ for $\lambda_1(A)$ with non-negative coordinates $v_i$. Set $t:=\min \{x_i/v_i: v_i\neq 0\}$ and note that for each coordinate $i$ we have $\lambda x_i\geq (A\mathbf{x})_i\geq \left(A(tv)\right)_i$ since $x_j\geq tv_j$ for all $j\in[n]$ and $A$ has nonnegative entries. It follows that $\lambda x_i\geq t(Av)_i=t\lambda_1(A)v_i$. Finally, there is a coordinate $j$ for which $x_j=tv_j\neq 0$ and plugging this into the inequality above yields $\lambda\geq \lambda_1(A)$. 
\end{proof}

From this point onward, the eigenvalues, eigenvectors, and spectrum of a graph will refer to those of its adjacency matrix, which is the square $(0,1)$-matrix that encodes the edges of the graph.

We conclude this section with a metric result which roughly says that if a set $X$ of bounded diameter is large enough, then the number of pairs of points from $X$ which are close to each other grows at least linearly with respect to $|X|$. We shall use this simple packing estimate to rule out the case where too few points lie in the boundary layer.

\begin{prop}\label{prop:many-close-pairs}
    For any real number $a\geq 1$ there is a constant $c_a>0$ such that for any $0<\eps<1$ and any finite set $X\subset \mathbb{R}^2$ of diameter at most $1$ with size $|X|\geq c_a\cdot \varepsilon^{-2}$, we have: 
    $$|\mathcal{N}_\eps(X)|\geq a|X|.$$
\end{prop}

\begin{proof}
    Since $X$ has diameter $1$, so does its convex hull $\text{conv}(X)$, meaning its area is bounded. Thus $\text{conv}(X)$ can be covered by a collection of $M:=\lfloor r\cdot\eps^{-2}\rfloor$ balls of radius $\eps/2$, where $r>0$ is a universal constant. We distribute the points of $X$ to these balls such that each point gets assigned to exactly one of the balls containing it. Suppose the $M$ balls contain $y_1,y_2,\ldots,y_M$ points from $X$ after the assignment, with $\sum_{i=1}^M y_i=|X|$. Since any pair of points assigned to the same ball is $\eps$-neighboring, we infer that ${\left|\mathcal{N}_\eps(X)\right|\geq \sum\limits_{i\in [M]}\dbinom{y_i}{2}= \dfrac{1}{2}\Bigl(\Bigl(\sum\limits_{i\in [M]}y_i^2\Bigr) - |X|\Bigr)}$.

    By the Cauchy-Schwarz inequality we get $\sum_{i=1}^M y_i^2\geq |X|^2/M$. Hence, for any $c_a\geq 2ra+r$, we obtain, as desired, that 
    
    $$\left|\mathcal{N}_\eps(X)\right|\geq\frac{|X|^2}{2M}-\dfrac{|X|}{2}\geq \frac{c_a\cdot\eps^{-2}\cdot |X|}{2r\cdot \eps^{-2}}-\dfrac{|X|}{2}\geq \frac{(c_a-r)\cdot|X|}{2r}\geq a|X|.$$
\end{proof}

\section{The box graph and the spectral reduction}\label{sec:reduction}

    In this section we define a natural auxiliary graph for a point set $X$ of finite diameter, whose vertices will represent clusters of neighboring points, while its edges will encode antipodal pairs. We then show how to reduce the problem to bounding the largest eigenvalue of this graph. 
    
    Let us start by observing that it is enough to prove Theorem \ref{thm:main} only for $\eps_1\leq \eps_2$. Indeed, if $\eps_1>\eps_2$, then any $\eps_2$-antipodal pair is $\eps_1$-antipodal as well. Hence, using the $\eps_2=\eps_1$ case, we obtain $\left|\mathcal{N}_{\eps_1}(X)\right|\geq  c_2\cdot \eps_1^{2}\cdot\eps_1^{-3/2}\left|\mathcal{A}_{\eps_1}(X)\right| = c_2\cdot \eps_1^{1/2}\left|\mathcal{A}_{\eps_1}(X)\right|\geq c_2\cdot \eps_1^{1/2}\left|\mathcal{A}_{\eps_2}(X)\right|$. Therefore, from now on we assume that $\eps_1\le\eps_2$.

%the result holds in this case, then $\left|\mathcal{N}_{\eps_1}(X)\right|\geq c_2\cdot \eps_1^{1/2}\left|\mathcal{A}_{\eps_1}(X)\right|\geq c_2\cdot \eps_1^{1/2}\left|\mathcal{A}_{\eps_2}(X)\right|$ when $\eps_1>\eps_2$. This happens since any pair $\{x_1,x_2\}\in \mathcal{A}_{\eps_2}(X)$ also contributes to $\mathcal{A}_{\eps_1}(X)$ for $\eps_1>\eps_2$.


%Let us start by observing that it is enough to prove Theorem \ref{thm:main} only for $\eps_1\leq \eps_2$. Indeed, if the result holds in this case, then $\left|\mathcal{N}_{\eps_1}(X)\right|\geq c_2\cdot \eps_1^{1/2}\left|\mathcal{A}_{\eps_1}(X)\right|\geq c_2\cdot \eps_1^{1/2}\left|\mathcal{A}_{\eps_2}(X)\right|$ when $\eps_1>\eps_2$. This happens since any pair $\{x_1,x_2\}\in \mathcal{A}_{\eps_2}(X)$ also contributes to $\mathcal{A}_{\eps_1}(X)$ for $\eps_1>\eps_2$.

Let $K:=\text{conv}(X)$ denote the convex hull of $X$. If a point $x\in X$ is at a distance greater than $\varepsilon_2$ from the boundary of $K$, i.e. if $\text{dist}(x,\partial K)>\varepsilon_2$, then $x$ does not contribute  any pair to $\mathcal{A}_{\eps_2}(X)$. Therefore, if we define $K_{\eps_2}:=\{x\in K:\text{dist}(x,\partial K)\le {\eps_2}\}$ and set $X_{\eps_2}=X\cap K_{\eps_2}$, we get $|\mathcal{A}_{\eps_2}(X_{\eps_2})|=|\mathcal{A}_{\eps_2}(X)|$, and $|\mathcal{N}_{\eps_1}(X_{\eps_2})|\leq |\mathcal{N}_{\eps_1}(X)|$. This shows that it is enough to prove the desired inequality for $X_{\eps_2}$ and hence from now on our focus will be on $X_{\eps_2}$. 

Recall that the perimeter of $K$ is bounded by a constant since it has diameter $1$, thus $K_{\eps_2}$, whose area is of order $\eps_2$, can be covered by $N=\Theta \left(\varepsilon_2\cdot\varepsilon_1^{-2}\right)$ squares of side length $\eps_1/2$ whose centers lie inside $K_{\eps_2}$ while the distance between any two of these centers is at least $\eps_1/4$. In other words, the center of a square is never covered by another square. We will refer to these squares as \emph{boxes} $B_1,B_2,\ldots,B_N$.%, where the labeling is made in cyclic order along the boundary $\partial K$, i.e. we pick a ray starting from the centroid of $K$ and we move it counterclockwise to label the boxes. \textcolor{red}{why centroid?}

The \emph{box graph} $G_X$ is the graph on vertex set $[N]:=\{1,2,3,\ldots,N\}$ where we connect $i$ and $j$ if there are $x\in X\cap B_i$ and $y\in X\cap B_j$ such that $\|x-y\|\geq 1-\eps_2$. We write $i\sim j$ to express that $i$ and $j$ are adjacent vertices in $G_X$ and also let $d_i:=\text{deg}_{G_X}(i)$. 
Furthermore, we assign each vertex of $X_{\eps_2}$ to one of the boxes containing it, let $n_i$ denote the number of points of $X$ that are assigned to the box $B_i$.

We note that the box graph is not unique for a set $X$, as there are different ways to cover $K_{\eps_2}$ by boxes of size $\eps_1/2$, but any one will be suitable for our purposes. We will use the notation $G_X(\eps_1,\eps_2)$ to make it clear that we refer to a box graph where we used squares of side length $\eps_1/2$ to cover $K_{\eps_2}$.  We will now see how $G_X(\eps_1,\eps_2)$ helps us compare $\mathcal{N}_{\eps_1}(X)$ with $\mathcal{A}_{\eps_2}(X)$.  Note that if $|\mathcal{N}_{\eps_1}(X)|\ge (\eps_1/\eps_2)^2 \cdot|\mathcal{A}_{\eps_2}(X)|$ then Theorem \ref{thm:main} holds, so it is enough to consider the case $|\mathcal{N}_{\eps_1}(X)|< (\eps_1/\eps_2)^2 \cdot|\mathcal{A}_{\eps_2}(X)|$.

\begin{prop}\label{prop:eigen}
    There is a universal constant $c_1>0$ such that for any $0<\eps_1\leq \eps_2<1$ the following statement holds. Suppose $X\subset \mathbb{R}^2$ is a finite set of diameter at most $1$, with size $|X|\geq c_1\varepsilon_1^{-2}$, for which $|\mathcal{N}_{\eps_1}(X)|< (\eps_1/\eps_2)^2 \cdot|\mathcal{A}_{\eps_2}(X)|$. Then for any box graph $G_X(\eps_1,\eps_2)$ associated to it, the following inequality holds: $$4\lambda_1\big(G_X(\eps_1,\eps_2)\big)\cdot|\mathcal{N}_{\eps_1}(X)|\geq |\mathcal{A}_{\eps_2}(X)|.$$
\end{prop}


\begin{proof}
 Suppose $G_X(\eps_1,\eps_2)$ is an $N$-vertex graph and let $B$ denote its adjacency matrix. Furthermore, let $\mathbf{n}:=(n_1,n_2,\ldots,n_N)$ be the vector whose $i^{\text{th}}$ coordinate counts how many points of $X_{\eps_2}$ are assigned to the box $B_i$.
    
On the one hand, we note that $|\mathcal{A}_{\eps_2}(X)|\leq \sum_{i\sim j}n_in_j=\langle\bf{n},\emph{B}\bf{n}\rangle$  since any pair $\{x,y\}\subset X$ with $\|x-y\|\geq 1-\eps_2$ must come from two boxes $B_i$ and $B_j$ which are neighbors in the box graph. If we show that $4|\mathcal{N}_{\eps_1}(X)|\geq \sum_{i=1}^Nn_i^2=\langle\bf{n},\bf{n}\rangle$, then we can conclude that:
    $$\dfrac{|\mathcal{A}_{\eps_2}(X)|}{|\mathcal{N}_{\eps_1}(X)|}\leq \dfrac{4\langle\bf{n},\emph{B}\bf{n}\rangle}{\langle\bf{n},\bf{n}\rangle}\leq 4\cdot \max_{\bf{x}\in\mathbb{R}^N\setminus\{0\}}\dfrac{\langle\bf{x},\emph{B}\bf{x}\rangle}{\langle\bf{x},\bf{x}\rangle}=4\lambda_1(B).$$


Hence, we focus on lower bounding $|\mathcal{N}_{\eps_1}(X)|$. Any pair of points that lie in the same box $B_i$ must be $\eps_1$-neighboring, thus  
$|\mathcal{N}_{\eps_1}(X)|\geq \displaystyle\sum_{i\in [N]}\dbinom{n_i}{2}= \dfrac{1}{2}\displaystyle\sum_{i\in [N]}n_i^2 - \dfrac{|X_{\eps_2}|}{2}$. We are only left to show that we can pick $c_1$ so that $\dfrac{|X_{\eps_2}|}{2}$ is smaller than $ \dfrac{1}{4}\displaystyle\sum_{i\in [N]}n_i^2$.
Indeed $\displaystyle\sum_{i\in[N]} n_i^2\geq N^{-1}|X_{\eps_2}|^2 $ by the Cauchy-Schwarz inequality. Hence  whenever $|X_{\eps_2}|\geq 4N$, we obtain  $\displaystyle\sum_{i\in[N]} n_i^2\geq\frac{|X_{\eps_2}|^2}{N}>4|X_{\eps_2}|$. 
In this case, it would follow that $|\mathcal{N}_\eps(X)|\geq \dfrac{1}{4}\displaystyle\sum_{i\in [N]}n_i^2$.


As $N=\Theta \left(\varepsilon_2\cdot\varepsilon_1^{-2}\right)$ there is a universal constant $c_3>4\eps_1^2\eps_2^{-1}N$ for any box graph. Using Proposition \ref{prop:many-close-pairs} with $a=c_3^2/2$, we can fix a universal constant $c_1\ge1$ such that  $|\mathcal{N}_{\eps_1}(X)|\geq (c_3^2/2)\cdot|X|$ whenever $|X|\geq c_1\cdot \eps_1^{-2}$. We claim that $|X_{\eps_2}|\geq 4N$ for this choice of $c_1$.

 Indeed, suppose $|X_{\eps_2}|<4N$. Then we get
$\left|\mathcal{A}_{\eps_2}(X)\right|\leq \dbinom{|X_{\eps_2}|}{2}<8N^2\le c_3^{2\ }\eps_1^{-4\ }\eps_2^{2}/2$. Also, using $|X|>c_1\eps_1^{-2}$ and $c_1>1$, we have  $|\mathcal{N}_{\eps_1}(X)|\geq (c_3^2/2)\cdot|X|\ge(c_3^2/2)\cdot\eps_1^{-2}$.
 This would give us $|\mathcal{N}_{\eps_1}(X)|\geq (\eps_1/\eps_2)^2\cdot|\mathcal{A}_{\eps_2}(X)|$, contradicting the assumption of the proposition. 
\end{proof}
 
% Finally, let $B$ denote the adjacency matrix of the box graph $G$ and let $\textbf{n}:=(n_1,n_2,\ldots,n_N)$. Our two inequalities may be rewritten as $|\mathcal{A}_\eps(X)|\leq \langle\bf{n},\emph{B}\bf{n}\rangle$ and $4|\mathcal{N}_\eps(X)|\geq \langle\bf{n},\bf{n}\rangle$. 


Therefore, our current focus is on finding an optimal upper bound for $\lambda_1\big(G_X(\eps_1,\eps_2)\big)$.
We conclude this section by observing the trivial bound given by $\lambda_1(G_X)\leq |V(G_X)|=\Theta\big(\varepsilon_2 \varepsilon_1^{-2}\big)$. Indeed, this follows from Proposition \ref{vector-test-criterion} applied with the vector $\mathbf{1}:=(1,1,\ldots,1)$.

\section{Geometry gets involved}\label{sec:geom}
We are yet to exploit the geometry of the box graph. For convenience, if $P_1$ and $P_2$ are two points in the plane we let $|P_1P_2|$ denote the length of the closed segment determined by them.

 We start with a simple lemma on convex quadrilaterals. The lemma itself is very similar in nature to Theorem \ref{thm:main}, it says that if we have two pairs of points at distance $1-\varepsilon$, then their endpoints cluster to some degree. 

\begin{lem}\label{quadrilateral}
    Let $0<\varepsilon<1$ and let $A,B,C,D$ be the vertices of a convex quadrilateral in counterclockwise order in the plane whose diameter is at most $1$. If the lengths of the segments $AD$ and $BC$ are $\geq 1-\varepsilon$, then $|AB|\cdot |CD|<3\varepsilon$. 
\end{lem}

\begin{proof}
     If either $|AB|< 3\eps$ or $|CD|< 3\eps$ then we are done. First, we claim that we can assume that $|AC|=|BD|=1$. Indeed, consider what happens if we rotate the point $C$ clockwise around $B$. Since $ABCD$ is convex, both $AC$ and $CD$ increase, and the quadrilateral remains convex until $C$ reaches the line $AB$, see Figure \ref{fig:rotate}. As $|AB|\ge 3\eps$, $|AB|+|BC|>1$ i.e., the length of $AC$ reaches $1$ before $C$ reaches the $AB$ line. Hence, we can assume that $|AC|=1$. Similarly, we can rotate $D$ clockwise around $A$ until $|DB|=1$. (The rotation might increase the length of $CD$ above 1, but this will not cause any problems in the following argument.)

     \begin{figure}[!ht]
         \centering
         \includegraphics[width=0.3\linewidth]{figures/figrotate.pdf}
         \caption{We can rotate $C$ around $B$ until $AC=1$, it only increases $|AC|$ and $|CD|$.}
         \label{fig:rotate}
     \end{figure}

     Suppose then $|AB|=2h$ with $3\varepsilon\leq 2h\leq 1$ and set a coordinate system in the plane such that $A=(-h,0)$ and $B=(h,0)$. Then $C$ and $D$ must lie on two circular arcs (see Figure \ref{fig:quadlemma}) and let $C=(c_x,c_y)$ and $D=(d_x,d_y)$ be their coordinates with $c_x,c_y,d_y\ge0$ and $d_x\le0$. We need to show that $|CD|<3\eps/2h$, or equivalently, that $(c_x-d_x)^2+(c_y-d_y)^2<9\eps^2\cdot(2h)^{-2}$. 

    \begin{figure}[!ht]
        \centering
        \includegraphics[width=0.4\linewidth]{figures/figquad.pdf}
        \caption{
        $C$ and $D$ must lie on the two small red arcs, respectively.}
        \label{fig:quadlemma}
    \end{figure}

   Since $|BC|\geq 1-\eps$ and $|AC|=1$, this translates to $(c_x-h)^2+c_y^2\geq (1-\eps)^2>1-2\eps$ and $(c_x+h)^2+c_y^2=1$, respectively. By subtracting the second inequality from the first we obtain $c_x\leq \eps/2h\leq 1/3$.  Analogously we get $-d_x\leq \eps/2h\leq 1/3$ by considering the point $D$. 

   To bound the $y$-axis coordinates $c_y$ and $d_y$ note that we can write $c_y=f(c_x)$ and $d_y=f(-d_x)$ for the map $f:[0,1/3]\to (0,1)$ given by $f(t):=\sqrt{1-(t+h)^2}$. By the mean value theorem,
   there is $z\in [0,1/3]$ such that $|c_y-d_y|=|c_x-(-d_x)|\cdot |f^\prime(z)|$, where $f^\prime:[0,1/3]\to(-\infty,0)$ is the derivative of $f$, given by $f^\prime(t):=-(t+h)/ \sqrt{1-(t+h)^2}$. It is not hard to see that the function $-f^\prime$ is increasing on $[0,1/3]$, hence $|f^\prime(z)|\leq |f^\prime(1/3)|\leq 5/\sqrt{11}<\sqrt{5}$. It follows that $(c_y-d_y)^2< 5(c_x+d_x)^2$.
    Using that $0<c_x\leq \eps/2h$ and $0<-d_x\leq \eps/2h$ we obtain $(c_x-d_x)^2<\eps^2/h^2$ and $(c_y-d_y)^2< 5(c_x+d_x)^2\le 5\eps^2/4h^2$. 
    
    It follows that $|CD|^2\le9\eps^2/4h^2$, which implies $|CD|\le 3\eps/2h=3\eps/|AB|$.
\end{proof}

We note that the sharp bound in Lemma \ref{quadrilateral} seems to be $|AB|\cdot|CD|\le 2\varepsilon-\eps^2$, however, any estimate of order $O(\eps)$ suffices for our purposes; therefore, we provided a weaker bound with a simpler proof. 

% %[
% \begin{figure}[!ht]
%         \centering
%         \includegraphics[width=0.3\linewidth]{figures/fig2.pdf}
%         \caption{Caption}
%         \label{fig:quadlemma}
%     \end{figure}

%     Straightforward calculation gives that the endpoints of the arcs are:

%     \[P_1,P_2=\left(
% \pm\frac{\varepsilon(2-\varepsilon)}{2d},\
%  \frac{\sqrt{\bigl((2-\varepsilon)^2-d^2\bigr)\bigl(d^2-\varepsilon^2\bigr)}}{2d}
% \right)
% \]

% \[Q=\left(0,\sqrt{1-\frac{d^2}{4}}\right)\]
% Therefore $C$ and $D$ differ by at most $\frac{\varepsilon(2-\varepsilon)}{d}$ in their $x$-coordinate. Next, we compare their $y$-coordinate.

% Let
% \[
% a:=1-\frac{d^2}{4},
% \qquad
% b:=\frac{\bigl((2-\varepsilon)^2-d^2\bigr)\bigl(d^2-\varepsilon^2\bigr)}{4d^2}.
% \]
% Then
% \[
% \sqrt{1-\frac{d^2}{4}}
% -
% \frac{\sqrt{\bigl((2-\varepsilon)^2-d^2\bigr)\bigl(d^2-\varepsilon^2\bigr)}}{2d}
% =
% \sqrt a-\sqrt b
% =
% \frac{a-b}{\sqrt a+\sqrt b}.
% \]
% A direct expansion yields
% \[
% a-b
% =
% \frac{\varepsilon(2-\varepsilon)\bigl(2d^2+2\varepsilon-\varepsilon^2\bigr)}{4d^2}.
% \]
% Since and \(2\varepsilon<d<1\), 
% \[
% 2d^2\le 2d,
% \qquad
% 2\varepsilon-\varepsilon^2\le 2\varepsilon<d,
% \]
% so
% \[
% 2d^2+2\varepsilon-\varepsilon^2<3d.
% \]
% Hence
% \[
% a-b
% \le
% \frac{\varepsilon(2-\varepsilon)\cdot 3d}{4d^2}
% \le
% \frac{3\varepsilon}{2d}.
% \]
% On the other hand,
% \[
% \sqrt a+\sqrt b\ge \sqrt a=\sqrt{1-\frac{d^2}{4}}
% \ge \sqrt{1-\frac14}
% =\frac{\sqrt3}{2},
% \]
% since \(d\le 1\). Therefore
% \[
% 0\le
% \sqrt a-\sqrt b
% =
% \frac{a-b}{\sqrt a+\sqrt b}
% \le
% \frac{\frac{3\varepsilon}{2d}}{\frac{\sqrt3}{2}}
% =
% \sqrt3\,\frac{\varepsilon}{d}
% <
% 2\frac{\varepsilon}{d}.
% \]
% Thus
% \[
% \sqrt{1-\frac{d^2}{4}}
% -
% \frac{\sqrt{\bigl((2-\varepsilon)^2-d^2\bigr)\bigl(d^2-\varepsilon^2\bigr)}}{2d}
% <
% 2\frac{\varepsilon}{d}.
% \]


Our next lemma  helps us identify, from a geometric perspective, the center of a box with its corresponding vertex in the box graph.   

\begin{lem}\label{box-point}
    Let $G:=G_X(\varepsilon_1,\varepsilon_2)$ be an $N$-vertex box graph associated to a finite set $X\subset \mathbb{R}^2$ of diameter at most $1$ and let $B_i$ and $B_j$ be the two boxes corresponding to vertices $i\sim j$ in $G$. Then $1-\eps_1-\eps_2< |P_iP_j|\leq 1$, where $P_i$ and $P_j$ are the centers of the boxes $B_i$ and $B_j$, respectively. 
\end{lem}

\begin{proof}
    Suppose the points $X_i$ and $X_j$ in $X$ certify that $i\sim j$, i.e. that $|X_iX_j|\geq 1-\eps_2$. We also know that $|P_iX_i|<\eps_1/2$ and $|P_jX_j|<\eps_1/2$, therefore by the triangle inequality we obtain that $|P_iP_j|>|X_iX_j|-|X_iP_i|-|X_jP_j|>1-\eps_2-\eps_1/2-\eps_1/2=1-\eps_1-\eps_2$. Furthermore, $P_i$ and $P_j$ lie inside $\text{conv}(X)$, hence $|P_iP_j|\leq 1$, finishing the proof.
\end{proof}

We are now ready to derive the key tail estimate on the local edge density of the box graph around a specific vertex.

\begin{lem}\label{lem:tails}
There is a universal constant $c_0>0$ for which the following statement holds. 
Let $0<\eps_1\leq \eps_2<1$ and let $G:=G_X(\eps_1,\eps_2)$ be an $N$-vertex box graph associated to a finite set $X\subset \mathbb{R}^2$ of diameter at most $1$. Then, for every vertex $k\in [N]$ and integer $t\ge 1$, we have:
\[
|\{j\sim k: d_j\ge t\}|\le c_0(N/t)\cdot (\eps_2/\eps_1)^2.
\]
\end{lem}

\begin{proof}
%As $N_t:=\{j\sim k: d_j\ge t\}$ is a subset of the neighborhood of $k$, its size is trivially at most $d_k$. It remains to prove that its size is also at most $c_0N/t\cdot (\eps_2/\eps_1)^2$ for a suitable constant $c_0>0$.
First note that if $c_0$ is large enough, then the statement holds for $t\le 5$. Hence, we may assume that $t>5$. Let $N_t=\{j\sim k: d_j\ge t\}$. Since $N=\Theta(\varepsilon_2\ \varepsilon_1^{-2})$, we can fix a constant $c_0$ such that
$c_0>2^{21}\cdot \eps_1^{-2}\ \eps_2N^{-1}$.
Suppose by contradiction that $|N_t|> c_0(N/t)\cdot (\eps_2/\eps_1)^2$. By the choice of $c_0$, we have $|N_t|> (2^{21}/t)\cdot (\eps_2^3/\eps_1^4)$.

Recall that we have labeled the boxes $B_1,\dots, B_N$ and let $P_1,\dots, P_N$ denote the centers of the boxes, respectively.  Let $<_k$ be the ordering of the elements of $N_t$ in counterclockwise order from the viewpoint of $P_k$, that is, $i<_kj$ if $P_i$ is to the left of the directed line $P_kP_j$.

\begin{figure}[!ht]
    \centering
    \includegraphics[width=0.5\linewidth]{figures/taillemma.pdf}
    \caption{Finding four boxes, $B_k,B_m,B_i,B_j$, that are far from each other.}
    \label{fig:tail}
\end{figure}

 Pick a middle box from $N_t$ according to $<_k$ and label it $B_m$. Clearly, it has at least $t$ neighbors. At least $(t-1)/2$ of these neighbors  correspond to boxes that lie in the same half-plane determined by the line passing through $P_k$ and $P_m$.
 Let the box $B_i$ be the one among these that is farthest from $P_k$. By the choice of box $B_m$, at least $\lfloor (|N_t|-1)/2\rfloor$ of the neighbors of $k$ correspond to boxes that lie in the half-plane determined by the line $P_kP_m$ not containing $P_i$, let $B_j$ be the farthest box in $N_t$ on that side, see Figure \ref{fig:tail}. 
 
 We claim that the four points $P_i,P_k,P_j,P_m$ contradict Lemma \ref{quadrilateral} for the parameter $2\eps_2$. That is, we claim that $|P_iP_m|>1-2\eps_2$, $|P_kP_j|>1-2\eps_2$ but $|P_iP_k|\cdot |P_jP_m|\ge 6\eps_2$. Using that $\eps_1\le \eps_2$, the inequalities $|P_iP_m|>1-2\eps_2$ and $|P_kP_j|>1-2\eps_2$ follow from Lemma \ref{box-point}. Hence, it only remains to bound $|P_iP_k|$ and $|P_jP_m|$.

First consider $|P_iP_k|$. Let $S$ be the annulus  sector that is centered at $P_m$, has inner and outer radii $1-\eps_2$ and $1$, respectively, and is bounded by the lines $P_mP_i$ and $P_m 
P_k$. By the choice of $P_i$ we know that at least $(t-1)/2$ boxes intersect $S$. As $t>5$ we have $(t-1)/2\ge t/3$. Recall that each box has side length $\eps_1/2$ and that their centers are $\eps_1/4$ apart. This implies that we can fit at least $t/3$ disjoint boxes of side length $\eps_1/4$ inside the annulus sector $S'$ centered at the same point $P_m$ and bounded by the lines $P_mP_i$ and $P_mP_k$, but with inner and outer radii $1-\eps_2-\eps_1/4$ and $1+\eps_1/4$, respectively. Hence, the area of $S'$ is at least $\eps_1^2\cdot t/48$. 
On the other hand, if we let $\alpha$ denote the measure of $\angle P_iP_mP_k$ in radians, then the area of $S'$ is also given by $\big((1+\eps_1/4)^2-(1-\eps_2-\eps_1/4)^2\big)\cdot\alpha/2\leq \big((1+\eps_2/4)^2-(1-5\eps_2/4)^2\big)\cdot\alpha/2\leq 3\eps_2\cdot \alpha/2$. However, $|P_iP_k|=2\sin(\alpha/2)\geq \alpha/2$ since $\alpha\leq \pi$, so the area of $S'$ is at most $|P_iP_k|\cdot 3\eps_2$. It follows from here that  $\eps_1^2\cdot t/48\le|P_iP_k|\cdot 3 \eps_2$, which implies that $|P_iP_k|\ge\eps_1^2\eps_2^{-1}\cdot t\cdot2^{-9}$.

 A similar argument gives $|P_jP_m|\ge\eps_1^2\eps_2^{-1}\cdot |N_t|\cdot 2^{-9}$, hence $|P_iP_k|\cdot |P_jP_m|\ge\eps_1^4\eps_2^{-2}\cdot t\cdot|N_t|\cdot 2^{-18}$. Using that $|N_t|> (2^{21}/t)\cdot (\eps_2^3/\eps_1^4)$, we obtain $|P_iP_k|\cdot |P_jP_m|\ge 8 \eps_2>6\eps_2$, as desired. We conclude that our assumption is false and so $|N_t|\leq c_0N/t\cdot (\eps_2/\eps_1)^2$.
\end{proof}

With these tools in hand, we are now ready to bound the largest eigenvalue of the box graph.

\begin{prop}\label{prop:sqrt-control}
    There exists a universal constant $c_0>0$ such that for any  $0<\eps_1\leq \eps_2<1$ and for any finite set of points $X\subset \mathbb{R}^2$ of diameter at most $1$, and any box graph $G:=G_X(\eps_1,\eps_2)$ associated to it, the following spectral bound holds:$$\lambda_1\big(G_X(\eps_1,\eps_2)\big)\leq 6\cdot\eps_{2 } \eps_1^{-1}\sqrt{c_0|V(G_X(\eps_1,\eps_2))|}.$$ 
\end{prop}

\begin{proof}
   Let $c_0$ be the constant given by Lemma \ref{lem:tails}.
   Let $V(G)=[N]$, $d_i=\text{deg}_{G_X}(i)$ for all $i\in[N]$ and let $\mathbf{r}:=(\sqrt{d_1},\sqrt{d_2},\ldots,\sqrt{d_N} )$. We claim that $(B\mathbf{r})_i\leq 6\eps_2\eps_1^{-1}\sqrt{c_0N\cdot d_i}$ for each $i\in[N]$, where $B$ is the adjacency matrix of $G$. This claim implies the desired result by Proposition \ref{vector-test-criterion}, finishing the proof.  To proceed, note that $(B\mathbf{r})_i=\sum_{j\sim i}\sqrt{d_j}$. To bound this sum, we will apply a standard dyadic interval decomposition combined with Lemma \ref{lem:tails}.
    
    For a fixed vertex $i$ we partition the neighborhood of $i$ into the following sets: first we place the neighbors of low degree into $A_0:=\{j\sim i:d_j\cdot d_i\leq c_0N\eps_2^{2\ }\eps_1^{-2}\}$; next, for each integer $k$ that ranges from $d:=\lfloor\log_2(c_0N\eps_2^{2\ }\eps_1^{-2}d_i^{-1})\rfloor$ to $\lfloor\log_2N\rfloor$, we set $A_k:=\{j\sim i:2^k\leq d_j<2^{k+1}\}$. From Lemma \ref{lem:tails} we can see that for such a $k$ the following holds: $$\sum_{j\in A_k}\sqrt{d_j}<\sqrt{2^{k+1}}\cdot|A_k|\leq \sqrt{2^{k+1}}\cdot|\{j\sim i:d_j\geq 2^k\}|\leq \sqrt{2^{k+1}}\cdot \dfrac{c_0N\eps_2^2}{2^{k\ }\eps_1^2}=\dfrac{c_0N\eps_2^2}{\eps_1^2\sqrt{2^{k-1}}}.$$

Moreover,  $\displaystyle\sum_{j\in A_0} \sqrt{d_j}\leq \eps_2\ \eps_1^{-1}\sqrt{c_0N/d_i}\cdot |A_0|\leq \eps_2\ \eps_1^{-1}\sqrt{c_0Nd_i}$, where we have used the trivial bound $|A_0|\leq d_i$. Putting all these together we get:
$$\sum_{j\sim i}\sqrt{d_j}\leq \dfrac{\eps_2}{\eps_1}\cdot\sqrt{c_0Nd_i}+\left(\dfrac{\eps_2}{\eps_1}\right)^2 \cdot\sum_{k=d}^{\lfloor \log_2N\rfloor}\dfrac{c_0N}{\sqrt{2^{k-1}}}\leq\dfrac{\eps_2}{\eps_1}\cdot \sqrt{c_0Nd_i} +\dfrac{c_0N\eps_2^2}{\eps_1^2\sqrt{2^{d-1}}}\cdot \sum_{k\geq 0}\left(\sqrt{2}\right)^{-k}.$$

However, $2^{d-1}>c_0N\eps_2^{2\ }\eps_1^{-2}(4d_i)^{-1}$ and $\sum_{k\geq 0}\alpha^k =(1-\alpha)^{-1}$ for each $0<\alpha<1$, therefore by plugging these in above we finally obtain the desired upper bound:
$$\sum_{j\sim i}\sqrt{d_j}\leq \dfrac{\eps_2}{\eps_1}\cdot \sqrt{c_0Nd_i}+\dfrac{2\eps_2}{\eps_1}\cdot\sqrt{c_0Nd_i}\cdot (1+\sqrt{2})\leq \dfrac{6\eps_2}{\eps_1}\cdot\sqrt{c_0N\cdot d_i}.$$
\end{proof}

Putting it all together, we obtain a proof of Theorem \ref{thm:main}.
\begin{proof}[Proof of Theorem \ref{thm:main}]
    As we have previously discussed, it is enough to prove the case ${\eps_1<\eps_2}$. If ${|\mathcal{N}_{\eps_1}(X)|/|\mathcal{A}_{\eps_2}(X)|\geq (\eps_1/\eps_2)^2}$ then we are done since  $(\eps_1/\eps_2)^2=\Omega(\eps_1^2\cdot \eps_2^{-3/2})$. On the other hand, if $|\mathcal{N}_{\eps_1}(X)|/|\mathcal{A}_{\eps_2}(X)|< (\eps_1/\eps_2)^2$, then by using Proposition \ref{prop:eigen} we get an absolute constant $c_0>0$ such that  $|\mathcal{N}_{\eps_1}(X)|/|\mathcal{A}_{\eps_2}(X)|\geq \left(4\lambda_1(G)\right)^{-1}$ for any $|X|\ge c_1\eps_1^{-2}$ and any box graph $G=G_X(\eps_1,\eps_2)$. Finally, Proposition \ref{prop:sqrt-control}, together with the fact that $|V(G)|=\Theta(\eps_2\ \eps_1^{-2})$, help us draw the conclusion that 
    $$\dfrac{|\mathcal{N}_{\eps_1}(X)|}{|\mathcal{A}_{\eps_2}(X)|}=\Omega\big(\lambda_1(G)^{-1}\big) =\Omega\left(\eps_1\eps_2^{-1}\sqrt{1/|V(G)|}\right)= \Omega\left( \eps_1^2\cdot \eps_2^{-3/2}\right).$$ 
\end{proof}


    
%With this last result, we have completed the proof of Theorem \ref{thm:main} by settling the case when $\eps_1\leq \eps_2$ as well. Indeed, as seen in Section \ref{sec:reduction}, we may assume that $|X_{\eps_2}|=\Omega(\eps_2\ \eps_1^{-2})$, otherwise we get the even stronger bound $|\mathcal{N}_{\eps_1}(X)|/|\mathcal{A}_{\eps_2}(X)|\geq (\eps_1/\eps_2)^2$. But then Proposition \ref{prop:eigen} applies to give us $|\mathcal{N}_{\eps_1}(X)|/|\mathcal{A}_{\eps_2}(X)|\geq \left(4\lambda_1(G)\right)^{-1}$, where $G:=G_X(\eps_1,\eps_2)$ is a box graph associated with $X$. Finally, Proposition \ref{prop:sqrt-control}, together with the fact that $|V(G)|=\Theta(\eps_2\ \eps_1^{-2})$, help us draw the conclusion that $|\mathcal{N}_{\eps_1}(X)|/|\mathcal{A}_{\eps_2}(X)|= \Omega\left( \eps_1^2\cdot \eps_2^{-3/2}\right)$. 
%\end{proof}

\section{Constructions}\label{sec:examples}

In this section, we provide several examples that show that the bounds from Theorem \ref{thm:main} are tight up to a constant factor. Moreover, we show that the key tail estimate from Lemma \ref{lem:tails}\linebreak is essentially sharp, thus capturing the extremal behavior of the box graph density.

\subsection{Annulus construction}
 A natural idea is to consider a set of points which are uniformly distributed near the boundary of a disk of radius $1/2$. Our first construction follows this idea, as $X$ will be a subset of an annulus with inner and outer radii $(1-\eps_2)/2$ and $1/2$, respectively. We simply place the points as concentric regular polygons, as in Figure \ref{fig:annulus}.

For convenience, given two integers $a\leq b$, we will denote by $\mathbb{Z}[a,b]$ the set $\{a,a+1,\ldots,b\}$, i.e. the set of integers ranging from $a$ to $b$. We will also write $\mathbb{Z}_n$ for the set of remainders modulo $n$, i.e., $\mathbb{Z}_n:=\mathbb{Z}[0,n-1]$ for any positive integer $n$.

\begin{figure}
    \centering
    \includegraphics[width=0.3\linewidth]{figures/constr1.pdf}
    \caption{The annulus construction for $r=20$, $s=4$.}
    \label{fig:annulus}
\end{figure}

\begin{exam}[Annulus construction]
    Let $0<\delta<1$ and let $r,s\geq 2$ be integers. We define:
    $$AU(\delta,r,s):=\left\{ \left(\dfrac{1}{2}-\dfrac{\delta l}{2s}\right)\left( \cos\left(\frac{\pi k}{r}\right), \sin\left(\frac{\pi k}{r}\right)\right) : (k,l)\in \mathbb{Z}[-r,r-1]\times \mathbb{Z}_s\right\}.$$
\end{exam}

Our goal is to show the following estimate.

\begin{prop}\label{ref:annularconstr}
    Let $0<2\eps<\delta\leq 1/2$ and let $r,s\geq \eps^{-2}$ be integers. Then:
    \begin{align*}
     \left|\mathcal{A}_{\delta}\big(AU(\delta,r,s)\big)\right|=&\ \Theta\left(\sqrt{\delta}\cdot r^2s^2\right); \\
     \left|\mathcal{N}_{\eps}\big(AU(\delta,r,s)\big)\right|=&\ \Theta\left((\eps\cdot rs)^2/\delta\right).
\end{align*}
\end{prop}

 This shows that if $r$ and $s$ are large enough, then the threshold in Theorem \ref{thm:main} is sharp in the $2\eps_1\leq \eps_2$ regime. Indeed, for $2\eps_1\leq \eps_2\le 1/2$ and $X=AU(\eps_2,r,s)$, it gives

 $$
\dfrac{|\mathcal{N}_{\eps_1}(X)|}{|\mathcal{A}_{\eps_2}(X)|} = \frac{\Theta\left((\eps_1\cdot rs)^2/\eps_2\right)}{\Theta\left(\sqrt{\eps_2}\cdot r^2s^2\right)} = \Theta(  \eps_1^2\cdot \eps_2^{-3/2}).$$


Before proving Proposition \ref{ref:annularconstr}, we introduce a useful lemma for counting the points in an annulus sector.

\begin{dfn}
    Given $-1\leq \alpha_1<\alpha_2<1$ and $0\leq \delta_1\leq \delta_2<1$, we define an annulus sector associated to these parameters as follows:
    $$AS(\alpha_1,\alpha_2,\delta_1,\delta_2):=\left\{ \frac{1-t}{2}\cdot \big( \cos\left(\pi \alpha\right), \sin\left(\pi\alpha\right)\big) : (\alpha,t)\in [\alpha_1,\alpha_2]\times [\delta_1,\delta_2]\right\}.$$
\end{dfn}

\begin{lem}\label{lem:annulus-count}
        Let $-1\leq \alpha_1<\alpha_2<1$, $0\leq \delta_1\leq \delta_2\leq \delta<1$ and $r,s\geq 2$ be integers. Then: 
        $$\big|AU(\delta,r,s)\cap AS(\alpha_1,\alpha_2,\delta_1,\delta_2)\big|=\big(1+\lfloor \alpha_2r\rfloor-\lceil \alpha_1r\rceil\big)\cdot\big(1+\lfloor \delta_2s/\delta\rfloor-\lceil  \delta_1s/\delta\rceil\big).$$
\end{lem}

\begin{proof}
 Consider the point $P\in AU(\delta,r,s)$ that has the form $(1/2-\delta l/2s)\big(\cos(\pi k/r),\sin(\pi k/r)\big)$ for some $(k,l)\in \mathbb{Z}[-r,r-1]\times \mathbb{Z}_s$. Then $P$ lies in $AS(\alpha_1,\alpha_2,\beta_1,\beta_2)$ if and only if its coordinates satisfy $\pi\alpha_1\leq \pi k/r\leq \pi\alpha_2$ and $(1-\delta_2)/2\leq 1/2-\delta l/2s\leq (1-\delta_1)/2$. However, after rearranging, this translates to $\alpha_1 r\leq k\leq \alpha_2 r$ and $\delta_1s/\delta\leq l\leq\delta_2s/\delta$, which implies the result. 
\end{proof}

%We are now ready to show that the annular construction is extremal for our main theorem, as computing the ratio between $\mathcal{N}_{\eps_1}(X)$ and $\mathcal{A}_{\eps_2}(X)$ for $0<2\eps_1\leq \eps_2\leq 1/2$ and $X=AU(\eps_2,r,s)$ would yield the desired $\eps_1^2\cdot \eps_2^{-3/2}$ threshold. \textcolor{red}{ r and s equals what?}


\begin{proof}[Proof of Proposition \ref{ref:annularconstr}]
Set $X:= AU(\delta,r,s)$ and start with estimating $\mathcal{N}_\eps(X)$. We claim that for every point $Q\in X$ the ball of radius $\eps$ centered at $Q$ contains $\Theta\left(\eps^2rs/\delta\right)$ other points of $X$. Since each of these points forms an $\eps$-neighboring pair with $Q$, and $|X|=2rs$, the claim does imply that $|\mathcal{N}_\eps(X)|=\Theta\left(\eps^2r^2s^2/\delta\right)$.

Due to the circular symmetry of our construction, we can assume the point $Q$ has coordinates $\big((s-\delta q)/2s,0\big)$ for some $q\in \mathbb{Z}_s$. Let $(k,l)\in  \mathbb{Z}[-r,r-1] \times \mathbb{Z}_s$ and let $P\in AU(\delta,r,s)$ be a point with coordinates
$(1/2-\delta l/2s)\big(\cos(\pi k/r),\sin(\pi k/r)\big)$. Then the distance $|PQ|$ satisfies: 
    \begin{align*}
        |PQ|^2 &= \left(\dfrac{1}{2}-\dfrac{\delta q}{2s}-\cos\left(\dfrac{\pi k}{r}\right)\cdot \left(\dfrac{1}{2}-\dfrac{\delta l}{2s}\right)\right)^2 +\left(\dfrac{1}{2}-\dfrac{\delta l}{2s}\right)^2\cdot \sin^2\left(\dfrac{\pi k}{r}\right) \\
         &=  \left(\dfrac{1}{2}-\dfrac{\delta q}{2s}\right)^2 - 2\cos\left(\dfrac{\pi k}{r}\right)\cdot \left(\dfrac{1}{2}-\dfrac{\delta q}{2s}\right)\left(\dfrac{1}{2}-\dfrac{\delta l}{2s}\right) + \left(\dfrac{1}{2}-\dfrac{\delta l}{2s}\right)^2\\
         &= \dfrac{\delta^2}{4s^2}\cdot (q-l)^2 + 2\left(1-\cos\left(\dfrac{\pi k}{r}\right)\right)\cdot \left(\dfrac{1}{2}-\dfrac{\delta q}{2s}\right)\left(\dfrac{1}{2}-\dfrac{\delta l}{2s}\right).
    \end{align*}
   \indent Let us note that $\dfrac{1}{2}\geq 2\left(\dfrac{1}{2}-\dfrac{\delta q}{2s}\right)\left(\dfrac{1}{2}-\dfrac{\delta l}{2s}\right)\geq\dfrac{1}{2}(1-\delta)^2=\dfrac{1}{8}$ since $1-\delta\geq \dfrac{1}{2}$. Moreover, we know that $x^2/4\leq 1-\cos x\leq x^2/2$ whenever $|x|<2.7$, which will be useful later as well. 
 
 Consequently, if $|k|\geq 2\varepsilon r$, then $|PQ|^2\geq (1-\cos(\pi k/r))/8\geq 2(\pi k/8r)^2>\varepsilon^2$. Similarly, if $|q-l|>2\varepsilon s/\delta$, then $|PQ|^2>\varepsilon^2$. This means that the ball of radius $\eps$ centered at $Q$ is contained in $AS(-2\eps r,\ 2\eps r,\ \delta q/s-\eps,\ \delta q/s +\eps)$. 
 
  On the other hand, if both $|q-l|\leq 6\eps s/5\delta$ and $|k|\leq \eps r/2$ hold, then we deduce from the above that $|PQ|^2\leq 9\eps^2/25 +\pi^2k^2/(2r)^2\leq 9\eps^2/25+16\eps^2/25=\eps^2$, which means that the ball of radius $\eps$ centered at $Q$ contains $AS(-\eps r/2,\ \eps r/2,\ \delta q/s-3\eps/5,\ \delta q/s+3\eps/5)$. It follows by Lemma \ref{lem:annulus-count} that any point $Q\in X$ contributes with $\Theta\left(\eps^2rs/\delta\right)$ pairs to $\mathcal{N}_{\eps}(X)$.

We now move to $\mathcal{A}_\delta(X)$. Similarly as before, we claim that for every point $Q\in X$ there are $\Theta\left(\sqrt{\delta}\cdot rs\right)$ points in $X$ which are $\delta$-antipodal with $Q$. Again, as $|X|=2rs$, this claim implies that $|\mathcal{A}_\delta(X)|=\Theta\left(\sqrt{\delta}\cdot r^2s^2\right)$.

We may assume this time that $Q$ has coordinates $\big((\delta q-s)/2s,0\big)$ for some $q\in \mathbb{Z}_s$, i.e. it is situated on the negative part of the $x$-axis due to the symmetry. We let $(k,l)\in \mathbb{Z}[-r,r-1]\times \mathbb{Z}_s$ and let $P\in AU(\delta,r,s)$ be the point with coordinates
$(1/2-\delta l/2s)\big(\cos(\pi k/r),\sin(\pi k/r)\big)$. We estimate $|PQ|$ in a similar fashion: 
 \begin{align*}
        |PQ|^2 &= \left(\dfrac{1}{2}-\dfrac{\delta q}{2s}+\cos\left(\dfrac{\pi k}{r}\right)\cdot \left(\dfrac{1}{2}-\dfrac{\delta l}{2s}\right)\right)^2 +\left(\dfrac{1}{2}-\dfrac{\delta l}{2s}\right)^2\cdot \sin^2\left(\dfrac{\pi k}{r}\right) \\
         &=  \left(\dfrac{1}{2}-\dfrac{\delta q}{2s}\right)^2 + 2\cos\left(\dfrac{\pi k}{r}\right)\cdot \left(\dfrac{1}{2}-\dfrac{\delta q}{2s}\right)\left(\dfrac{1}{2}-\dfrac{\delta l}{2s}\right) + \left(\dfrac{1}{2}-\dfrac{\delta l}{2s}\right)^2\\
         &= \left(1-\delta\cdot \dfrac{q+l}{2s}\right)^2  - 2\left(1-\cos\left(\dfrac{\pi k}{r}\right)\right)\cdot \left(\dfrac{1}{2}-\dfrac{\delta q}{2s}\right)\left(\dfrac{1}{2}-\dfrac{\delta l}{2s}\right).
    \end{align*}
  \indent Recalling the useful bounds from above, we deduce that $|PQ|^2\leq 1-\big(1-\cos(\pi k/r)\big)/8\leq 1-\pi^2k^2/32r^2$. Therefore, if $|k|\geq 8r\sqrt{\delta}/\pi$, then $|PQ|^2< 1-2\delta<(1-\delta)^2$, which implies that all points of $X$ which form $\delta$-antipodal pairs with $Q$ must be in $AS(-8r\sqrt{\delta}/\pi,8r\sqrt{\delta}/\pi,0,\delta)$.  

 Finally, we claim that if $2l\leq s$ and $|k|\leq \sqrt{2\delta}\cdot r/\pi$, then $|PQ|^2\geq (1-\delta)^2$, which would imply that all points of $X$ from inside $AS(-\sqrt{2\delta}\cdot r/\pi,\sqrt{2\delta}\cdot r/\pi,0,\delta/2)$ are $\delta$-antipodal with $Q$. Consequently, Lemma \ref{lem:annulus-count} applied to these two annulus sectors would then give us that there are $\Theta\left(\sqrt{\delta}\cdot rs\right)$ such points, as desired.

 To prove this last claim, we have $\left(1-\delta\cdot \dfrac{q+l}{2s}\right)^2\ge 4\left(\dfrac{1}{2}-\dfrac{\delta q}{2s}\right)\left(\dfrac{1}{2}-\dfrac{\delta l}{2s}\right)$ by AM-GM,
 while also $2\big(1-\cos(\pi k/r)\big)\leq \pi^2k^2/r^2$. It follows that
 $$|PQ|^2\geq \left(4-\dfrac{\pi^2k^2}{r^2}\right)\left(\dfrac{1}{2}-\dfrac{\delta q}{2s}\right)\left(\dfrac{1}{2}-\dfrac{\delta l}{2s}\right)\geq \left(1-\dfrac{\pi^2k^2}{4r^2}\right)\left(1-\dfrac{\delta}{2}\right)(1-\delta).$$
  Therefore, it is enough to prove that $1-(\pi k/2r)^2\geq (1-\delta)(1-\delta/2)^{-1}=1-\delta(2-\delta)^{-1}$. But this is true since $(\pi k/2r)^2\leq (\sqrt{2\delta}/2)^2=\delta/2<\delta(2-\delta)^{-1}$, finishing the proof.
\end{proof}


\subsection{Reuleaux triangle constructions}

In this subsection we study an extremal construction which is based on the Reuleaux triangle. Suppose $P_0,P_1,P_2$ are points in $\mathbb{R}^2$ forming an equilateral triangle of side length $1$. We draw, for each $i\in \{0,1,2\}$, an arc between $P_i$ and $P_{i+1}$ centered at $P_{i+2}$, where the indices are understood modulo $3$. The resulting shape is called the \emph{Reuleaux triangle}. It has a constant width of 1 across all directions and a diameter of 1. 


\begin{exam}[Reuleaux construction]
    Let $0<\eps<1$ and let points $P_0,P_1,P_2$ be the corners of a Reuleaux triangle of diameter $1$. For each $i\in \{0,1,2\}$, consider the points $A_i, B_i,C_i,D_i$ on the arc $P_{i+1}P_{i+2}$ such that $|A_iP_{i+1}|=|B_iP_{i+2}|=1/3$ and $|C_iP_{i+1}|=|D_iP_{i+2}|=\eps/2$, where all the indices are modulo $3$. We define the set $RT(\eps,x,y)$ in the following way: for each $i\in\{0,1,2\}$ we place $y$ equally spaced points along each of the arcs $A_iB_i$, then we place $x$ points inside the triangle $\mathcal{T}_i:=P_iC_{i-1}D_{i+1}$.
    (See Figure \ref{fig:reucon1})  
\end{exam}



For large enough values of $x$ and $y$, we can sharply determine $\mathcal{N}_{\eps_1}(X)$ and $\mathcal{A}_{\eps_2}(X)$ for the set $X:=RT(\eps_2,x,y)$ whenever $2\eps_1\geq \eps_2$, as illustrated by the next result.

\begin{prop}\label{2nd-example}
    Let $0<\eps_2\leq2\eps_1<3^{-3}$ and let $x,y$ be large enough integers. Then:
    \begin{align*}
     \left|\mathcal{A}_{\eps_2}\big(RT(\eps_2,x,y)\big)\right|=&\ 3x^2+3xy; \\
     \left|\mathcal{N}_{\eps_1}\big(RT(\eps_2,x,y)\big)\right|=&\ \Theta\big(\eps_1\cdot y^2+ x^2\big).
%      \begin{cases}
% \ \ \ \ \ \ \ \Theta\big(\eps_1\cdot y^2+ x^2\big), & 2\eps_1 \geq \eps_2, \\
% \Theta\left(\eps_1\cdot y^2+\left(\dfrac{\eps_1}{\eps_2}\right)^2\cdot \dbinom{x}{2}\right),& 2\eps_1<\eps_2.
% \end{cases}
\end{align*}
\end{prop}

Setting $y:=tx$ for some parameter $t\geq 1$ we get $\left|\mathcal{N}_{\eps_1}\big(RT(\eps_2,x,y)\big)\right|=\Theta\left(x^2(1+t^2\eps_1)\right)$ and $\left|\mathcal{A}_{\eps_2}\big(RT(\eps_2,x,y)\big)\right|=\Theta (tx^2)$ above. By picking $t=\Theta\left(\eps_1^{-1/2}\right)$, it is now easy to observe that we achieve the $\sqrt{\eps_1}$ threshold from Theorem \ref{thm:main} for the ratio $|\mathcal{N}_{\eps_1}(X)|/|\mathcal{A}_{\eps_2}(X)|$.

\begin{proof}[Proof of Proposition \ref{2nd-example}.]
    Set $X:=RT(\eps_2,x,y)$ and  start with $\mathcal{A}_{\eps_2}(X)$. We claim that all pairs from $\mathcal{A}_{\eps_2}(X)$ must have one of their points in some $\mathcal{T}_i$. Indeed, notice that the longest distance that can occur between two points of $X$ lying on the three arcs $A_iB_i$ is given by $|A_iB_{i+1}|$. These three segments are congruent, so let us consider $|A_1B_2|$. For the convex quadrilateral $A_1B_2P_1P_2$, we have $|A_1P_2|\cdot |B_2P_1|=3^{-2}>3\eps_2$, hence we infer from Lemma \ref{quadrilateral} that $|A_1B_2|<1-\eps_2$. 


\begin{figure}[!ht]
    \centering
    \includegraphics[width=0.6\linewidth]{figures/reu1.pdf}
    \caption{The Reuleaux triangle and the first construction based on it.}
    \label{fig:reucon1}
\end{figure}
    
    Next, let us count how many $\eps_2$-antipodal pairs a point $Q\in X\cap \mathcal{T}_1$ forms in $X$. 
    We start by noting that the longest distance from $Q$ to any point of $X$ lying on one of the arcs $A_2B_2$ and $A_3B_3$ is given by $|QA_{2}|$ and $|QB_{3}|$, respectively. However, the angle $\angle QA_{2}P_{3}$ is obtuse, hence $|QA_{2}|^2<|QP_{3}|^2-|A_{2}P_{3}|^2\leq 1-3\eps_2\leq 1-2\eps_2-\eps_2^2$, which gives $|QA_{2}|<1-\eps_2$. Similarly, 
    $|QB_{3}|<1-\eps_2$, implying that only the points from inside $\mathcal{T}_{2},\mathcal{T}_{3}$ and those on the arc $A_1B_1$ form $\eps_2$-antipodal pairs with $Q$. We claim that all of them do. Indeed, for any point $C$ on the arc $A_1B_1$ we have $|CQ|\geq|P_1C|-|P_1Q|\geq 1-\eps_2/2$, while for any point $D$ inside $\mathcal{T}_j$ with $j\neq 1$ we have $|QD|\geq |P_iP_j|-|P_iQ|-|P_jD|\geq 1-\eps_2$. Hence $Q$ appears in $2x+y$ antipodal pairs. The same argument works for $Q\in X\cap \big(\mathcal{T}_2\cup \mathcal{T}_3\big)$. In total we obtain $|\mathcal{A}_{\eps_2}(X)|=3x^2+3xy$.   
    
   Finally, we shift our focus to $\mathcal{N}_{\eps_1}(X)$. Clearly, if two points are $\eps_1$-neighboring, then they must lie inside the same triangle $\mathcal{T}_i$ or on the same arc $A_iB_i$. Note that if $y$ is large enough, then each point from one of the arcs $A_iB_i$ forms $\Theta(\eps_1y)$ $\eps_1$-neighboring pairs within $X$. Hence $\Theta(\eps_1y^2)$ neighboring pairs come from the arcs. Moreover, any two points inside the same $\mathcal{T}_i$ are $\eps_1$-neighboring and thus $3x(x-1)/2$ other $\eps_1$-neighboring pairs get added to $\mathcal{N}_{\eps_1}(X)$. This yields the desired bounds on $|\mathcal{N}_{\eps_1}(X)|$, finishing the proof.
\end{proof}

\medskip

We conclude this subsection with another construction based on the Reuleaux triangle. Its purpose is to show that the degree bound for the box graph from Lemma~\ref{lem:tails}, which is a central ingredient to the proof, is tight up to a constant factor.

\begin{exam}\label{exa:sharptail}

    Let $0<\delta<1$ and let points $P_0,P_1,P_2$ be the corners of an equilateral triangle of side length $1$ whose center is $O$.
    Consider the points $Q_0,Q_1,Q_2$ in the plane such that $Q_i$ lies on the segment $OP_i$ such that $|P_iQ_i|=\delta/\sqrt{3}$, thus $Q_0,Q_1,Q_2$ form an equilateral triangle of side length $1-\delta$. Draw the two Reuleaux triangles determined by the triples $\{P_i\}$ and $\{Q_i\}$ and define the set $RA(\delta)$ to be the region between the two Reuleaux triangles (see Figure \ref{fig:reucon2}).
\end{exam}

\begin{figure}[!ht]
    \centering
    \includegraphics[width=0.3\linewidth]{figures/reu2fixed.pdf}
    \caption{The region $RA(\delta)$.}
    \label{fig:reucon2}
\end{figure}

\begin{prop}\label{sharp-deg-tail}
    For all $0<2\varepsilon_1\leq \varepsilon_2<0.5$ there is a covering of the region $RA\big(\eps_2/\sqrt{3}\big)$ whose associated box graph $G:=G(\eps_1,\eps_2)$  has at least $2\eps_2/\eps_1$ vertices $v$ which, for any positive integer $t$, satisfy the following degree bound:
    $$\big|\{w\in V(G): w\sim v \text{ in } G \text{ and deg}(w)\geq t \}\big|\geq \left\lfloor \left (\dfrac{\eps_2}{\eps_1}\right)^2 \cdot \dfrac{|V(G)|}{4t}\right\rfloor\ .$$
\end{prop}

\begin{proof}
    Take a covering of $RA(\eps_2/\sqrt{3})$ by boxes of size $\eps_1/2$ where no two box centers are closer to each other than $\eps_1/4$. Let $P_0P_1P_2$ and $Q_0Q_1Q_2$ denote the outer and the inner equilateral triangles of $RA(\eps_2/\sqrt{3})$. The arc $P_iP_{i+1}$ has length $\pi/3$ and the radial segment $P_iQ_i$ has length $\eps_2/3$. A simple calculation shows that we can associate a box graph on $N=\big(4.5+o(1)\big) \cdot \eps_{2\ }\eps_1^{-2}$ vertices to $RA(\eps_2/\sqrt{3})$, where, as in Section \ref{sec:reduction}, two boxes are connected if we can find a point in each of them such that the two form an $\eps_2$-antipodal pair.

    Due to the symmetry, it is enough to prove the bound for a vertex $v$ whose box $B_v$ intersects the segment $P_1Q_1$, since there are $2\eps_2/3\eps_1$ of them. Fix $t\geq 1$ and consider the points $A_2$ and $A_1$ on the arcs $Q_1Q_0$ and $Q_2Q_0$, respectively, such that $|Q_1A_2|=t\eps_1^{2\ }\eps_2^{-1}$ and $|A_1A_2|=1-\eps_2$.\linebreak Moreover, let $Q'_1,A_2'$ denote the intersection points of the arc $P_0P_1$ with the lines $P_2Q_1$ and $P_2A_2$, respectively. Similarly, let $Q'_2,A'_1$ denote the intersection points of the arcs $P_0P_2$ with the lines $P_1Q_2$ and $P_1A_1$, respectively (see Figure \ref{fig:reuc2b}).

    \begin{figure}
        \centering
        \includegraphics[width=0.3\linewidth]{figures/reu2.pdf}
        \caption{$\mathcal{S}_1$ and $\mathcal{S}_2$ span a complete bipartite graph in the box graph.}
        \label{fig:reuc2b}
    \end{figure}
    
    Next, we note that the shortest distance between a point inside the annulus sector $Q_1A_2A'_2Q'_1$ and a point inside the annulus sector $Q_2Q'_2A'_1A_1$ is given by the segment $A_1A_2$ of length $1-\eps_2$. 
    
    This means the subgraph of $G$ induced by the boxes covering these two regions is a complete bipartite graph. Let $\mathcal{S}_1$ and $\mathcal{S}_2$ denote its two vertex sets. Also note that the points in $\mathcal{S}_2$ are at distance at least $1-\eps_2$ from $v$. Our goal is to show that $|\mathcal{S}_1|\geq t$ and $|\mathcal{S}_2|\geq (\eps_2/2\eps_1)^2\cdot N/t$. This would finish the proof since each vertex $w\in \mathcal{S}_2$ has all vertices from $\mathcal{S}_1$ in their neighborhood, which implies that $\mathcal{S}_2\subset \{w\sim v: \text{deg}(w)\geq t\}$, giving  $|\{w\sim v: \text{deg}(w)\geq t\}|\geq (\eps_2/2\eps_1)^2\cdot N/t$.

   To bound $|\mathcal{S}_1|$, let $\alpha_1$ denote the measure of angle $\angle Q_1P_2A_2$ in radians and recall the area of the sector $Q_1A_2A'_2Q'_1$ is given by $\big(1-(1-\varepsilon_2/3)^2\big)\cdot \alpha_1/2>\alpha_1(1-\eps_2/3)\cdot \eps_2/3\geq 4t/3\cdot (\eps_1/2)^2$, where the last inequality follows as $\alpha_1(1-\eps_2/3)$ actually represents the length of the arc $Q_1A_2$, which is bigger than $|Q_1A_2|=t\eps_1^{2\ }\eps_2^{-1}$.
    Therefore, we must have at least $t$ vertices of $G$ that correspond to boxes inside the sector $Q_1A_2A'_2Q'_1$. This means that $|\mathcal{S}_1|\geq t$ and, if we show that $|A_1Q_2|\geq 2\eps_{2}N/9t$, the same argument will show that the sector $Q_2Q'_2A'_1A_1$ has area bigger than $(2/3)^3\cdot (\eps_2/\eps_1)^2\cdot N/t\cdot \eps_1^2/4$, which implies that $|\mathcal{S}_2|\geq (\eps_2/2\eps_1)^2\cdot N/t$. Thus, we are only left with bounding $|A_1Q_2|$.    

   By applying Ptolemy's inequality in the quadrilateral $Q_1A_2A_1Q_2$ with $|A_1A_2|=1-\eps_2$ and $|Q_1Q_2|=|Q_1A_1|=|Q_2A_2|=1-\eps_2/\sqrt{3}$, we get $|A_1Q_2|\cdot |A_2Q_1|\geq |A_1Q_1|\cdot|A_2Q_2|-|A_1A_2|\cdot|Q_1Q_2|$, leading to $|A_1Q_2|\cdot |A_2Q_1|\geq \left(1-\eps_2/\sqrt{3}\right)\left(\eps_2-\eps_2/\sqrt{3}\right)\geq \eps_2\left(1-1/\sqrt{12}\right)\left(1-1/\sqrt{3}\right) >\eps_2/4$. However, this can be rewritten as $|A_1Q_2|>\eps_1^{-2\ }\eps_2^2\cdot (4t)^{-1}\ge2\eps_{2}N/9t$, as desired.
\end{proof}

\section{Concluding remarks and open questions}\label{sec:conclude}

The box graph arises naturally when one seeks to analyze the relationship between antipodal and neighboring pairs. After reducing the problem to estimating the largest eigenvalue of its adjacency matrix, one of the main difficulties was the selection of an appropriate test vector for the Collatz–Wielandt bound in Proposition~\ref{vector-test-criterion}.
Using the all-ones vector, it is easy to get an upper bound of $N^{3/2}$. However, a box graph may contain vertices of very large degree, potentially linear in $N$, which renders the all-ones vector unsuitable. A more refined candidate is the degree vector $\mathbf{d}$, since $A\mathbf{d}$   counts the number of $2$-paths originating at each vertex, where $A$ denotes the adjacency matrix of $G$. Nonetheless, this choice only yields an  $\sqrt{N\log N}$ upper bound for the order of the largest eigenvalue, leading to the $\sqrt{\varepsilon^{-1}\log(\varepsilon^{-1})}$ threshold, which was also obtained by  Korsky \cite{korsky2026optimalspectralboundsantipodal}.
Moreover, constructions such as in Example \ref{exa:sharptail} demonstrate that this selection cannot give a better bound. However, in all such examples, the vertices which generate a large number of $2$-paths appear to be relatively sparse. 
This suggests averaging over longer paths as a way to mitigate this concentration phenomenon. Indeed, our calculations suggested that by looking at $A^2\mathbf{d}$ --  which corresponds to counting $3$-paths -- we might be able to get a sharp result. Finally, we arrived at the vector $(\sqrt{d_1},\sqrt{d_2},\ldots,\sqrt{d_N} )$ by noting that it works for all examples known to us, including Example \ref{exa:sharptail}. \medskip
    


The problem is also worth investigating in higher dimensions. As illustrated by the annulus construction (Example 1), the extremal behavior in $\mathbb{R}^2$ of the ratio between $\mathcal{N}_{\eps_1}(X)$ and $\mathcal{A}_{\eps_2}(X)$ in the regime $\varepsilon_1 < \varepsilon_2$ can be achieved by comparing the area of a ball of radius $\varepsilon_1/2$ with that of the subset of a ball of radius $1/2$ consisting of points at distance at least $1-\varepsilon_2$ from a fixed point located near the boundary. In contrast, when $\varepsilon_1 > \varepsilon_2$, the corresponding threshold appears to depend solely on $\varepsilon_1$. We conjecture that an analogous phenomenon persists in higher dimensions: for $d \geq 3$, one of the extremal configurations should be realized by a set of points distributed uniformly in a region $\mathcal{R}$ of $\mathbb{R}^d$ defined as the space between two concentric balls $B_1$ and $B_2$ of diameters $1$ and $1 - \varepsilon_2$, respectively. 

A ball of radius $\eps_1/2$ has volume of order $\Theta_d\big(\eps_1^d\big)$, and if we fix a point $Q$ in $\mathcal{R}$, it can be shown that the subset of $\mathcal{R}$ made of points which are $\eps_2$-antipodal with $Q$ has volume $\Theta_d\big(\eps_2^{(d+1)/2}\big)$. When $\eps_1<\eps_2$, plenty of balls of radius $\eps_1/2$ can be fitted inside $\mathcal{R}$, leading to the volume ratio threshold. When $\eps_1>\eps_2$ this is no longer the case, and the behavior is the same as for $\eps_1=\eps_2$, since extremal constructions can be pushed closer to the boundary of $B_1$. 

Motivated by these observations, we propose the following conjecture.

\begin{conjecture}
    There are universal constants $c_1,c_2>0$ such that for any $0<\varepsilon_1,\varepsilon_2<1$ and any finite set $X\subset \mathbb{R}^d$ of diameter at most $1$ with $|X|\ge c_1\varepsilon_1^{-d}$ the following holds: 
\[
\dfrac{|\mathcal{N}_{\eps_1}(X)|}{|\mathcal{A}_{\eps_2}(X)|} \ge
\begin{cases}
c_2\cdot \eps_1^d\cdot \eps_2^{-(d+1)/2}, & \eps_1 \le \eps_2; \\
 c_2\cdot \eps_1^{(d-1)/2}, & \eps_1 > \eps_2.
\end{cases}
\]
\end{conjecture}

The reduction to determining the largest eigenvalue of the box graph naturally carries over to higher dimensions. Thus, extending our approach to $\mathbb{R}^d$ amounts to finding a suitable analog of Lemma \ref{quadrilateral} and understanding the corresponding tail bounds it would imply on the degree distribution of the box graph.

Another natural direction is to study the problem in different metric spaces. Already in $\mathbb{R}^2$, the answer depends on the choice of norm. 
For instance, consider the $\ell_\infty$ norm, and place $n/2$ points evenly along the segment from $(0,0)$ to $(0,1)$, while the remaining $n/2$ points are evenly placed along the segment from $(1,0)$ to $(1,1)$. Then there are $n^2/4$ pairs at distance exactly $1$ and $\Theta(n^2 \varepsilon)$ $\varepsilon$-neighboring pairs. This shows the lower bound of order $\varepsilon^{1/2}$ is not attainable in this setting.


\section*{Acknowledgment}
The authors wish to thank J\'anos Pach and Zolt\'an L\'or\'ant Nagy for the fruitful discussions about the question. The use of Lemma \ref{lem:tails} was suggested by ChatGPT 5.4, but the proof of the lemma is due to the authors.  

G\'abor Dam\'asdi was supported by the ERC Advanced Grant no. 882971 GeoScape, by the J\'anos Bolyai Research Scholarship of the Hungarian Academy of Sciences, the STARTING-25 project no. 152199 and the EXCELLENCE-24 project no. 151504 Combinatorics and Geometry of the NRDI Fund.

Lauren\c{t}iu Ploscaru was supported by the ERC Advanced Grants no. 882971 GeoScape and ERMiD no. 101054936.

\printbibliography

\end{document}
